% =====================================================================
% Chapter 4 --- The Perception Pipeline (cohesion hub)            [cross-cutting]
% Presents the end-to-end pipeline ONCE and maps the three self-contained
% project chapters (5, 6, 7) onto it. This is the chapter that keeps the
% three projects reading as one coherent work.
% Target length: ~5 pages.
% =====================================================================

\section{Overview: A Modular Camera--Point-Cloud Pipeline}
% TODO: end-to-end diagram RGB -> 2D detection -> frustum -> PointNet 3D -> EKF tracker -> 3D track;
% the pipeline as the realization of the "Sense" capability for detect-and-avoid.

\section{The Four Stages at a Glance}
\subsection{2D Detection}
\subsection{Frustum Proposal}
\subsection{3D Detection (PointNet)}
\subsection{Tracking (EKF)}
% TODO: one concise conceptual paragraph per stage (inputs/method/outputs); the detailed
% methodology is deferred to the project chapters, to avoid duplication.

\section{Modularity and Cross-Domain Generality}
% TODO: why a modular pipeline -- each stage is trained/validated independently and can be
% reused or swapped across domains; this is what makes the three projects composable.

\section{Mapping the Contributions onto the Pipeline}
% TODO: locate the three self-contained projects on the pipeline (the unifying thread):
%   - the EKF tracker stage          -> Chapter~\ref{cha:tracking}  (simulation; global vs. local)
%   - the monocular real-world variant -> Chapter~\ref{cha:monocular} (real UAV data; TU Berlin)
%   - the detection front-end + dataset -> Chapter~\ref{cha:airsim}    (synthetic data; AirSim)
% State explicitly: one vision (the pipeline), three contributions.
